\documentclass[12pt, a4paper]{report}

\usepackage[utf8]{inputenc}
\usepackage[spanish]{babel}
\usepackage{graphicx}
\usepackage{geometry}
\usepackage{hyperref}
\usepackage{listings}
\usepackage{color}
\usepackage{amsmath}
\usepackage{booktabs}
\usepackage{longtable}

\geometry{left=25mm, right=25mm, top=25mm, bottom=25mm}

\definecolor{codegreen}{rgb}{0,0.6,0}
\definecolor{codegray}{rgb}{0.5,0.5,0.5}
\definecolor{codepurple}{rgb}{0.58,0,0.82}
\definecolor{backcolour}{rgb}{0.95,0.95,0.92}

\lstdefinestyle{mystyle}{
    backgroundcolor=\color{backcolour},
    commentstyle=\color{codegreen},
    keywordstyle=\color{magenta},
    numberstyle=\tiny\color{codegray},
    stringstyle=\color{codepurple},
    basicstyle=\ttfamily\footnotesize,
    breakatwhitespace=false,
    breaklines=true,
    captionpos=b,
    keepspaces=true,
    numbers=left,
    numbersep=5pt,
    showspaces=false,
    showstringspaces=false,
    showtabs=false,
    tabsize=2
}
\lstset{style=mystyle}

\title{\textbf{RJ2XCL: Sistema Multilenguaje de Cómputo Científico\\Integrado en Microsoft Excel\\con Visor WebView2 y Notebooks Pluto.jl}}
\author{Minor Bonilla G.\\Universidad de Costa Rica\\Tesis de Maestría}
\date{19 de abril de 2026}

\begin{document}

\maketitle

\begin{abstract}
RJ2XCL es un sistema de cómputo científico multilenguaje que integra R y Julia como motores de cálculo dentro de Microsoft Excel, utilizando una arquitectura de procesos aislados comunicados mediante Named Pipes y Protocol Buffers. El sistema expone más de 90 procedimientos estadísticos y matemáticos accesibles directamente desde celdas de Excel, un visor HTML embebido basado en WebView2 para visualizaciones interactivas, y un modo avanzado con notebooks Pluto.jl para usuarios técnicos. La validación se sustenta en una suite de 200 pruebas unitarias y de integración con 0\% de fallos, incluyendo pruebas basadas en propiedades (Property-Based Testing). El sistema evolucionó de una calificación de calidad de 4.3/10 a 8.9/10 mediante correcciones sistemáticas de seguridad, confiabilidad y mantenibilidad.
\end{abstract}

\tableofcontents
\newpage

% ═══════════════════════════════════════════════════════════════════════════
\chapter{Introducción}
% ═══════════════════════════════════════════════════════════════════════════

\section{Contexto y Motivación}

En las ciencias de datos modernas existe una brecha entre el poder del procesamiento estadístico avanzado (R, Julia) y la accesibilidad en la capa de negocio (Microsoft Excel). Las soluciones existentes como BERT (Basic Excel R Toolkit) adolecen de problemas de estabilidad al inyectar librerías dinámicas directamente en el proceso de Excel, donde una violación de segmento del motor analítico provoca el colapso del proceso anfitrión.

RJ2XCL resuelve este problema mediante una arquitectura de tres capas con aislamiento de procesos: el XLL (add-in de Excel) se comunica con procesos hijo independientes (ControlR.exe, ControlJulia.exe) a través de Named Pipes, garantizando que un fallo en R o Julia no afecte la estabilidad de Excel.

\section{Objetivos}

\begin{enumerate}
    \item Aislar la ejecución del código analítico en procesos separados del entorno UI de Excel.
    \item Gestionar determinísticamente la memoria intercambiada con Excel mediante RAII en C++.
    \item Exponer funciones estadísticas y matemáticas de R y Julia como funciones nativas de Excel.
    \item Integrar un visor HTML embebido (WebView2) para visualizaciones interactivas.
    \item Proporcionar un modo avanzado con notebooks Pluto.jl para usuarios técnicos.
    \item Garantizar la reproducibilidad de análisis mediante exportación a notebooks.
\end{enumerate}

\section{Contribuciones}

Las contribuciones principales de este trabajo son:
\begin{itemize}
    \item Una arquitectura IPC robusta basada en Named Pipes y Protocol Buffers que soporta múltiples lenguajes de programación.
    \item La corrección de bugs críticos en la C API de Julia 1.12 (\texttt{jl\_arrayset} y \texttt{jl\_array\_ptr} para tipos primitivos inline).
    \item Un sandbox de seguridad que bloquea más de 90 patrones de ejecución peligrosa en R y Julia.
    \item La integración de WebView2 como visor HTML embebido en Excel con comunicación bidireccional JavaScript$\leftrightarrow$C++.
    \item Un generador de presentaciones HTML autocontenidas basado en reveal.js.
\end{itemize}

% ═══════════════════════════════════════════════════════════════════════════
\chapter{Arquitectura del Sistema}
% ═══════════════════════════════════════════════════════════════════════════

\section{Estructura de Tres Capas}

RJ2XCL implementa una topología de tres capas:

\begin{enumerate}
    \item \textbf{Capa de Presentación (XLL):} Add-in de Excel compilado como DLL en C++17 con MSVC 2022. Registra funciones UDF, gestiona el visor WebView2, y coordina la comunicación con los motores.
    \item \textbf{Capa de Transporte (IPC):} Named Pipes de Windows con serialización Protocol Buffers v3.21.12. Comunicación síncrona bidireccional con timeouts configurables por lenguaje.
    \item \textbf{Capa de Motores:} Procesos hijo independientes --- ControlR.exe (R 4.4.1) y ControlJulia.exe (Julia 1.12.6) --- que ejecutan código y retornan resultados serializados.
\end{enumerate}

\section{Componentes Principales}

\begin{table}[h]
\centering
\begin{tabular}{lll}
\toprule
\textbf{Componente} & \textbf{Tecnología} & \textbf{Función} \\
\midrule
XLL (RJ2XCL64.xll) & C++17, MSVC 2022 & Add-in de Excel, registro de funciones \\
ControlR.exe & C++ + R 4.4.1 & Motor R con embedding via R API \\
ControlJulia.exe & C++ + Julia 1.12.6 & Motor Julia con embedding via C API \\
ViewerManager & C++ + WebView2 SDK & Visor HTML embebido \\
PlutoManager & C++ & Gestión del servidor Pluto.jl \\
PresentationBuilder & C++ & Generador de presentaciones reveal.js \\
NotebookLibrary & C++ & Gestión de 13 notebooks precargados \\
NotebookExporter & C++ & Exportación de análisis reproducibles \\
SandboxVerifier & C++ & Validación de seguridad de código \\
ConfigService & C++ + JSON & Configuración centralizada \\
\bottomrule
\end{tabular}
\caption{Componentes principales del sistema RJ2XCL}
\end{table}

\section{Protocol Buffers para IPC}

El intercambio de datos utiliza Protocol Buffers con un esquema extensible:

\begin{lstlisting}[language=C++, caption={Esquema Protobuf simplificado}]
message Variable {
  oneof value {
    int32 integer = 5;
    double real = 6;
    string str = 7;
    bool boolean = 8;
    Array arr = 10;
    HtmlContent html_content = 16; // WebView2
  }
  string name = 15;
}
\end{lstlisting}

El campo \texttt{HtmlContent} (field 16) fue agregado para transportar contenido HTML desde los motores de lenguaje al visor WebView2.

\section{Gestión de Memoria: Patrón RAII}

La interfaz \texttt{XLOPER12} de Excel requiere gestión manual de memoria. RJ2XCL encapsula toda interacción en la clase \texttt{RaiiXlOper}, garantizando liberación determinista:

\begin{lstlisting}[language=C++, caption={Patrón RAII para XLOPER12}]
{
    RaiiXlOper result;
    Excel12(xlUDF, &result, 1, arg1);
    // Al salir del bloque, el destructor libera
    // la memoria de XLOPER12 automaticamente.
}
\end{lstlisting}

% ═══════════════════════════════════════════════════════════════════════════
\chapter{Funcionalidades Implementadas}
% ═══════════════════════════════════════════════════════════════════════════

\section{Funciones Estadísticas en R}

Se implementaron 8 módulos R con más de 60 procedimientos estadísticos, accesibles como \texttt{=R.func(datos, TipoOutput)}:

\begin{table}[h]
\centering
\begin{tabular}{llc}
\toprule
\textbf{Módulo} & \textbf{Paquete R} & \textbf{Procedimientos} \\
\midrule
Regresión Lineal & stats, stargazer & 12 \\
Regresión Binaria & stats & 8 \\
Datos de Panel & plm & 15 \\
Series de Tiempo & tseries, mFilter & 12 \\
Pronóstico & forecast & 10 \\
Efectos Mixtos & lme4 & 8 \\
Supervivencia & survival & 7 \\
Psicometría & psych & 7 \\
Bayesiana & rstanarm & 8 \\
Estadística Base & stats & 12 \\
Descriptiva & Hmisc & 10 \\
Supuestos & car, lmtest & 12 \\
\bottomrule
\end{tabular}
\caption{Módulos estadísticos R implementados}
\end{table}

\section{Funciones Científicas en Julia}

Se implementaron 9 funciones Julia con 61 procedimientos, accesibles como \texttt{=J.func(datos, TipoOutput)}:

\begin{table}[h]
\centering
\begin{tabular}{llc}
\toprule
\textbf{Función} & \textbf{Área} & \textbf{Procedimientos} \\
\midrule
JM\_Algebra & Álgebra Lineal & 12 \\
JM\_Calculo & Cálculo Numérico & 7 \\
JM\_EDO & Ecuaciones Diferenciales & 4 \\
JML\_Clasificacion & Machine Learning & 6 \\
JML\_Clustering & K-Medias & 6 \\
JML\_Estadistica & Estadística & 8 \\
JO\_Optimizar & Optimización & 7 \\
JC\_Transformar & Transformación de Datos & 6 \\
JC\_Utilidades & Utilidades & 5 \\
\bottomrule
\end{tabular}
\caption{Funciones Julia implementadas}
\end{table}

\section{Visor WebView2 Embebido}

El sistema integra Microsoft Edge WebView2 como visor HTML embebido dentro de Excel, permitiendo renderizar visualizaciones interactivas (Plotly, D3.js, htmlwidgets) sin abrir un navegador externo. El visor soporta:

\begin{itemize}
    \item Ventanas flotantes modeless asociadas a la ventana de Excel.
    \item Comunicación bidireccional JavaScript$\leftrightarrow$C++ via PostMessage.
    \item Sandbox de seguridad con navegación restringida.
    \item Múltiples instancias simultáneas con evicción FIFO.
\end{itemize}

\section{Modo Avanzado: Pluto.jl}

Para usuarios técnicos, el sistema ofrece un modo avanzado basado en Pluto.jl:

\begin{itemize}
    \item Servidor Pluto.jl lanzado como proceso Julia independiente.
    \item 13 notebooks precargados (7 R via RCall.jl, 5 Julia nativo, 1 mixto).
    \item Exportación de análisis como notebooks reproducibles.
    \item RCall.jl para pipelines mixtos Julia+R.
\end{itemize}

\section{Presentaciones reveal.js}

El sistema genera presentaciones HTML autocontenidas usando reveal.js:

\begin{lstlisting}[language=C++, caption={Flujo de creación de presentaciones desde Excel}]
=RJ2XCL.PRESENTATION.NEW("Resultados Q1")
=RJ2XCL.PRESENTATION.ADD.SLIDE("presentation-1",
    "# Resumen\nVentas crecieron 15%", "text")
=RJ2XCL.PRESENTATION.BUILD("presentation-1", "")
\end{lstlisting}

% ═══════════════════════════════════════════════════════════════════════════
\chapter{Seguridad}
% ═══════════════════════════════════════════════════════════════════════════

\section{Sandbox de Ejecución}

El \texttt{SandboxVerifier} bloquea más de 90 patrones de ejecución peligrosa:

\begin{itemize}
    \item \textbf{R:} \texttt{system()}, \texttt{shell()}, \texttt{eval(parse())}, \texttt{do.call()}, \texttt{.Call()}, \texttt{url()}, \texttt{Sys.setenv()}
    \item \textbf{Julia:} \texttt{ccall}, \texttt{@ccall}, \texttt{cglobal}, \texttt{unsafe\_*}, \texttt{include()}, backtick literals
    \item \textbf{Bypass prevention:} Detección de concatenación de strings, whitespace stripping, \texttt{getattr} patterns
\end{itemize}

\section{Seguridad del Visor WebView2}

\begin{itemize}
    \item Navegación restringida: solo \texttt{file://} dentro del directorio de datos y \texttt{about:blank}.
    \item DevTools deshabilitado por defecto (configurable para debugging).
    \item JavaScript habilitado para interactividad (Plotly, D3.js) pero sin acceso a filesystem o red.
    \item Extensión de política para \texttt{localhost:port} cuando Pluto está activo.
\end{itemize}

% ═══════════════════════════════════════════════════════════════════════════
\chapter{Correcciones Técnicas Significativas}
% ═══════════════════════════════════════════════════════════════════════════

\section{Bugs de Julia 1.12: Array I/O}

Julia 1.12 cambió el almacenamiento de arrays de tipos primitivos (\texttt{Float64}, \texttt{Int64}) a \textit{inline storage}. Las funciones \texttt{jl\_arrayset} y \texttt{jl\_array\_ptr} dejaron de funcionar correctamente para estos tipos:

\begin{lstlisting}[language=C++, caption={Fix para arrays primitivos en Julia 1.12}]
// Antes (roto en Julia 1.12):
jl_arrayset(julia_array, element, i);

// Despues (correcto):
if (array_base_type == jl_float64_type) {
    jl_array_data(julia_array, double)[i] = 
        jl_unbox_float64(element);
}
\end{lstlisting}

Adicionalmente, \texttt{jl\_array\_data} cambió de firma: \texttt{jl\_array\_data(a)} $\rightarrow$ \texttt{jl\_array\_data(a, type)}.

\section{Base.Docs.doc() en Julia 1.12}

La función \texttt{Base.Docs.doc()} lanza \texttt{MethodError} en Julia 1.12 para funciones sin docstrings formales. Se resolvió con \texttt{try/catch} en \texttt{ListFunctions()}.

% ═══════════════════════════════════════════════════════════════════════════
\chapter{Validación y Pruebas}
% ═══════════════════════════════════════════════════════════════════════════

\section{Suite de Pruebas}

\begin{table}[h]
\centering
\begin{tabular}{lcc}
\toprule
\textbf{Categoría} & \textbf{Tests} & \textbf{Estado} \\
\midrule
Sandbox (R + Julia + Python) & 109 & 100\% \\
Property-Based (Python sandbox) & 3 & 100\% \\
Confiabilidad (unit + PBT) & 35 & 100\% \\
Config, Security, Discovery & 20 & 100\% \\
Type conversions, RAII, COM & 12 & 100\% \\
Integration, Callbacks, Basic & 9 & 100\% \\
Misc (Result, Version) & 12 & 100\% \\
\midrule
\textbf{Total} & \textbf{200} & \textbf{100\%} \\
\bottomrule
\end{tabular}
\caption{Suite de pruebas automatizadas}
\end{table}

\section{Evolución de Calidad}

\begin{table}[h]
\centering
\begin{tabular}{lccc}
\toprule
\textbf{Hito} & \textbf{Nota} & \textbf{Tests} & \textbf{Cambio Clave} \\
\midrule
Estado original & 4.3 & 39 & Prototipo con deuda técnica \\
Post-correcciones & 6.8 & 119 & Seguridad, RAII, mutex \\
Post-Julia sysimage & 7.2 & 119 & Julia arranca en segundos \\
Post-Mantenibilidad & 7.9 & 165 & Constants, logging, Doxygen \\
Post-Confiabilidad & 8.1 & 200 & Health monitoring, errores \\
Post-Julia Libraries & 8.4 & 200 & 61 procedimientos Julia \\
Post-WebView2+Pluto & 8.9 & 200 & Visor, notebooks, presentaciones \\
\bottomrule
\end{tabular}
\caption{Evolución de la calificación de calidad}
\end{table}

% ═══════════════════════════════════════════════════════════════════════════
\chapter{Conclusiones y Trabajo Futuro}
% ═══════════════════════════════════════════════════════════════════════════

\section{Conclusiones}

RJ2XCL demuestra que es viable construir un sistema multilenguaje robusto que integre R y Julia en Microsoft Excel mediante una arquitectura de procesos aislados. Las contribuciones principales incluyen:

\begin{enumerate}
    \item Una arquitectura IPC que soporta múltiples lenguajes sin comprometer la estabilidad de Excel.
    \item Más de 120 procedimientos estadísticos y matemáticos accesibles como funciones nativas de Excel.
    \item Un visor HTML embebido con WebView2 para visualizaciones interactivas.
    \item Un modo avanzado con Pluto.jl para reproducibilidad y personalización.
    \item Una suite de 200 pruebas con 100\% de cobertura en paths críticos.
\end{enumerate}

\section{Trabajo Futuro}

\begin{itemize}
    \item \textbf{Callback thread:} Re-habilitar el hilo de callbacks para inserción directa de gráficos en Excel.
    \item \textbf{CI/CD:} Pipeline de integración continua con tests automatizados en cada commit.
    \item \textbf{Property-Based Testing:} Implementar las 20 propiedades formales definidas en el diseño del visor WebView2.
    \item \textbf{Python:} Resolver los problemas de estabilidad de ControlPython.exe para re-habilitar Python como tercer lenguaje.
\end{itemize}

% ═══════════════════════════════════════════════════════════════════════════
\chapter*{Agradecimientos}
% ═══════════════════════════════════════════════════════════════════════════

A la Universidad de Costa Rica por el apoyo institucional. Al equipo de desarrollo de BERT (Basic Excel R Toolkit) cuyo código base sirvió como punto de partida. A la comunidad de Julia por el soporte en la migración a Julia 1.12.

\end{document}
